\scanpage{71}
\uline{Variante} laissons tomber l'hypothèse $d'\ (=\mathrm{pgcd}\text{ des deg. des 0-cycles sur }X_\eta)=1$.
On sait que $\operatorname{Br}(K)$ est un grpe de $p$-torsion, donc on a
\[
0\longrightarrow\operatorname{Pic}(X_\eta)\mathbin{-}P(K)\mathbin{-}\Theta\mathbin{-}0,
\]
où $\Theta\subset\operatorname{Br}(K)$ est un grpe de $p$-torsion, d'où résulte que ${}_n\operatorname{Pic}(X_\eta)\simeq{}_nP(K)$. Donc il reste vrai que $u_n$ est isom. à ${}_nv$. Mais $\Gamma$ peut avoir de la torsion. On trouve
\diagram{\begin{tikzcd}[column sep=small]
0\arrow[r]&{}_n\Gamma\arrow[r]\arrow[d,"\unsure{\simeq}"']&\mathrm H^1(X,\mu_n)\arrow[r,dash]&\mathrm H^1(X_{\bar\eta},\mu_n)^I\arrow[r]&\Gamma_n\arrow[r,dash]&\operatorname{NS}(X_0)_n\arrow[r]&{}\\
&\mathbf Z/(n,d)\mathbf Z
\end{tikzcd}}
\deleted{\unsure{Soit}} et par passage à la limite sur les puissances de $\ell$
\diagram{\begin{tikzcd}[column sep=small]
0\arrow[r,dash]&\mathrm H^1(X,\mathbf Z_\ell(1))\arrow[r]&\mathrm H^1(X_{\bar\eta},\mathbf Z_\ell(1))^I\arrow[r]\arrow[dr]&\Gamma\otimes\mathbf Z_\ell\arrow[r,dash]&\operatorname{NS}(X_0)\otimes\mathbf Z_\ell\\
&&&R\otimes\mathbf Z_\ell\arrow[r,dash]&0
\end{tikzcd}}
$R=\operatorname{Ker}(\Gamma\mathbin{-}\operatorname{NS}(X_0))$.
\margin{N.B. $R\otimes\mathbf Z_\ell\simeq\begin{cases}\mathbf Z/\ell\mathbf Z&\text{si }\ell\mid d\\0&\text{sinon}\end{cases}$}
On sait aussi (tout au moins modulo rés. des singularités, \unsure{qu'on pourrait faire cas. de ce bon}) que le noyau $R$ de $\Gamma\mathbin{-}\operatorname{NS}(X_0)$ est de \uline{torsion}. \deleted{(donc un grpe cyclique \unsure{d'ordre diviseur de} $d$, $\unsure{\mathbf Z/d'\mathbf Z}$)} C'est donc un sous-grpe de $\operatorname{tors}\Gamma\simeq\mathbf Z/d\mathbf Z$, [avec une nette tendance d'être égal à $\mathbf Z/d\mathbf Z$ tout entier ; \deleted{\unsure{et}} tout au moins le noyau contient-il la partie de $\operatorname{tors}\Gamma$ première à $p$, \unsure{pt} argument de Raynaud utilisant un rev. étale convenable de $S$ $\ldots$) On a donc, si
\margin{\unsure{même sans résolution} cf \unsure{in} Raynaud \unsure{note}, \unsure{prop. 7 (ii)}}
\margin{demander Raynaud si $\operatorname{Ker}$ est donc $\mathbf Z/d\mathbf Z$}
\[
\Gamma'=\Gamma/T\deleted{\unsure{orsion}},\quad\deleted{\unsure{\&}}\quad T:\operatorname{Torsion}\text{ de }\Gamma
\]
\[
0\mathbin{-}T_n\mathbin{-}\Gamma_n\mathbin{-}\Gamma'_n\longrightarrow0
\]
d'où, choisissant un splitting de $\Gamma\mathbin{-}\Gamma'\oplus T$, on trouve
\[
\operatorname{Ker}(\Gamma_n\mathbin{-}\operatorname{NS}(X_0)_n)\simeq T_n+\underbrace{\operatorname{Ker}(\Gamma'_n\mathbin{-}\operatorname{NS}(X_0)_n)}_{\operatorname{Ker}(\overline\Gamma_n\mathbin{-}\operatorname{NS}(X_0)_n)}
\]
où $\overline\Gamma\subset\operatorname{NS}(X_0)$
\[
\overline\Gamma=\operatorname{Im}\Gamma.
\]
On trouve donc, comme $\overline\Gamma$ n'a plus que de la $p$-torsion
\diagram{\begin{tikzcd}[column sep=large]
0\arrow[r,dash]&\delta'\arrow[r]\arrow[d,equal]&\delta\arrow[r,dash]&\delta''\arrow[r,dash]&0\\
&\text{tors de }\operatorname{NS}(X_0)\text{ première à }p
\end{tikzcd}}
\margin{70}
